\chapter{Relative Rigidity of Lattice Handles}
\label{ch:relative-rigidity}

This chapter records the relative PL statement used for the three lower handle
types.  The adjective relative refers to the complete cylinder frontier, the
marked rims of the cutting surfaces, and the compact core protected by the
preceding Wall and prime-replacement constructions.  All three are part of the
conclusion needed for gluing.

\section{The marked lattice model}
\label{sec:relative-rigidity-model}
\label{sec:relative-rigidity-contract}

Let $\iota$ and $\kappa$ be finite index types, let
$L\subseteq\mathbb R^\kappa$ be a discrete rank-$|\kappa|$ lattice, and put
\[
 X=(\mathbb R^\iota)\times(\mathbb R^\kappa/L),\qquad
 H=\overline B_{\infty}^{\,\iota}(0,1)\times(\mathbb R^\kappa/L).
\]
The ball is the closed unit ball for the sup norm on the finite function
space.  Its marked frontier is
\[
 \partial H=\{u: \|u\|_\infty=1\}\times(\mathbb R^\kappa/L).
\]
The ambient space is boundaryless; $H$ is the displayed closed subset.  A
standard lattice atlas is a PL-domain atlas whose affine charts have the exact
inverse formula
\[
 d_j^{-1}(z)=\bigl((a_jz)_1,\,[ (a_jz)_2 ]_L\bigr)
\]
for an affine equivalence $a_j:\mathbb R^3\simeq
\mathbb R^\iota\times\mathbb R^\kappa$, on the chart target.

Here PL-irreducibility means that the domain is a PL domain and each
embedded chartwise PL sphere in its interior bounds a chartwise PL ball
there.

\begin{theorem}[Marked relative rigidity]
\label{def:marked-relative-torus-rigidity}
For two chart families $e,d$, if $|\iota|+|\kappa|=3$, $|\iota|\leq2$, $d$
is standard, both families are PL-irreducible on $H$, and
\[
 \phi:H\longrightarrow H,\qquad \phi^{-1}(\partial H)=\partial H,\qquad
 \phi\simeq\operatorname{id}_H\ \operatorname{rel}\ \partial H
\]
is continuous and chartwise PL from $e$ to $d$, there is a homeomorphism
$g:H\to H$ such that $g$ is chartwise PL from $e$ to $d$,
$g^{-1}$ is chartwise PL from $d$ to $e$, and
\[
 g\simeq\phi\ \operatorname{rel}\ \partial H,\qquad
 g\simeq\operatorname{id}_H\ \operatorname{rel}\ \partial H.
\]
\end{theorem}
The two homotopies are distinct data: the first replaces the supplied map and
the second permits a relative lift. They are homotopies of maps, not
isotopies. An ambient isotopy will be constructed separately below.
We prove the three cases after establishing the geometric inputs.

\section{The two geometric producers}
\label{sec:relative-rigidity-producers}
\label{sec:constructed-inputs}

\subsection{Brown's ball and descent to the quotient}

Let $S\subset\mathbb R^3$ be a locally flat embedded two-sphere.  Local pair
charts give a bicollar and two complementary components.  The bounded
component has compact closure $D$ and full frontier $S$. Work in the
one-point compactification. Collapse the two regions beyond the ends of
the bicollar to opposite poles, mapping its middle sphere to the equator.
These two regions are the only nonsingleton fibers. Brown's two-fiber
theorem makes them cellular; cancellation of a cellular fiber on each
side identifies the closures of the complementary domains with balls.
These are the cellular-fiber and complementary-domain results of
\cite[Theorems 1--5, pp.~74--76]{Brown1960};
the bicollar follows from \cite[Theorem 3, p.~339]{Brown1962}.
The construction retains both the sphere and the chosen closed side.
In particular, restricting the resulting compactification homeomorphism
uses two separate inverse-image equations, for the equator and the finite
closed ball. It gives a
homeomorphism of pairs
\[
 (D,S)\cong(\overline B^3,S^2).
\]
The boundary correspondence here is setwise. An arbitrary initial
parametrization of $S$ is not claimed to be the restriction of this particular
homeomorphism.

For a PL sphere in the interior of a lattice handle, lift it to the
covering product and apply a
linear equivalence with $\mathbb R^3$.  Brown gives a compact lifted ball
$D_0$. Let $U$ be the inverse image of $\operatorname{int}D_0$ under that
linear equivalence. Its closure is compact, its interior and frontier are
connected, and its frontier is the lifted sphere. The additive quotient
$\pi$ is injective on this frontier. For $0\ne v\in\ker\pi$, the frontiers of
$U$ and $v+U$ are disjoint. Connectedness and separation imply that their
closures are disjoint or one contains the other. Nesting is impossible:
choose a linear functional $\ell$ with $\ell(v)>0$ and maximize it on
$\overline U$. The inclusion $v+\overline U\subset\overline U$ contradicts
that maximum; the reverse inclusion gives the same contradiction with
$-v$. Thus the closures are disjoint. A collision of $\pi$ on $\overline U$
would give precisely such a nonzero translate, so $\pi$ is injective there.

The compact-to-Hausdorff restriction is a homeomorphism onto its image.
Local openness of $\pi$ identifies the entire image frontier with $S$.
Finally every interval coordinate is strictly between $-1$ and $1$ on
the compact lifted sphere. A nonconstant affine function cannot have a
maximum in the open bounded region, so the same strict bounds hold on
its closure. The descended ball consequently lies in the handle interior.

\subsection{Boundary-proper relative approximation}

Let $R$ be compact, with PL-domain structures $e,d$. Suppose a relatively open set
$U$ contains the whole frontier and the identity is chartwise PL on $U$.
The complement $R\setminus U$ is compact in the interior.  The source first
constructs a compact polyhedral carrier $C$ with
\[
 R\setminus U\subset\operatorname{int}C,\qquad
 \operatorname{fr}C\subset U,\qquad C\subset\operatorname{int}R.
\]
More precisely, choose a finite complex $K$ and a homeomorphism
$F:C\to|K|$ PL in the source atlas. On each chart piece an affine
projection from the finite realization recovers that source chart.
The inverse $F^{-1}:|K|\to R$ is already PL in the target atlas near
$F(\operatorname{fr}C)$, because the two structures agree on $U$.
Choose a closed subpolyhedron containing a neighborhood of this set and
contained in that agreement region. Hudson's relative approximation
\cite[Lemma 4.2 and homotopy-control remark, pp.~92--94]{Hudson1969}
gives a target-PL map $q$ and a homotopy from $F^{-1}$ to $q$ fixed on
the retained subpolyhedron. Choose its error small enough to keep the
whole homotopy in $\operatorname{int}R$.

On $C$ compose with $F$; on $R\setminus\operatorname{int}C$ use the
identity. These closed pieces agree on the entire $\operatorname{fr}C$
throughout, so they paste to a continuous homotopy $S$. The fixed neighborhood gives an
open overlap: the endpoint is PL on $\operatorname{int}C$ as $qF$, and
on the exterior neighborhood as the original identity. The construction
keeps the exact equation $S_t^{-1}(\partial R)=\partial R$ for every time
$t$, since the moving part stays in the interior.

At the endpoint, the map $\varphi:R\to R$ is chartwise PL, is homotopic to
the identity relative to $\partial R$, and satisfies the exact equation
\[
 \varphi^{-1}(\partial R)=\partial R.
\]
The reverse inclusion comes from the fixed collar; the forward inclusion comes
from the fact that the modified part remains in $\operatorname{int}R$.
The empty-domain case is the identity map.  This is the project-specific
boundary-proper adapter for Hudson's relative approximation theorem.

\section{Finite ordinary projection and marked cuts}
\label{sec:relative-rigidity-cuts}
\label{sec:marked-hierarchy}

An annulus embedded in one stage of the two-sheeted tower need not remain
embedded on projection to the preceding stage. We need an embedded annulus
whose two rims remain in their separately specified frontier marks
$F_0,F_1$, and remain essential in the original region. The marks are
disjoint and relatively open in that frontier. The tower and its finite
proper-surface constructions are those of
Section~\ref{sec:dehn-annulus-descent}.

First put the upper annulus in relative finite PL position. Intersections
of its finitely many chart triangles give a finite one-dimensional double
relation; this relation is generally not a finite set of points. Only the
exceptional pairs, where the projected branches do not yet have ordinary
crossing charts, form a finite set. At an interior exceptional value
flatten one branch, subdivide the other, and move its free vertices
slightly off the finitely many forbidden affine levels. Extend the
motion over the surrounding coordinate ball, fixing its outer collar.
At a frontier value do the same in a half-space chart: the zero-height
subcomplex is exactly the clipped rim, so tangential rim motions and
interior motions preserve the half-space and the prescribed frontier
mark. Common subdivision makes the motions PL and gives the full
two-branch image equations, including at old edges.

Choose disjoint closed windows about the finitely many exceptional
values, and perform the repairs in their inverse-image windows. Their
finite composite preserves the upper embedding. Inside a changed window
its whole branch images agree with those of that single repair; outside
the closed union of the change supports they agree with the old ordinary
branches. The windows cover every old exceptional value. This proves
ordinary crossings everywhere after projection, with no triple points,
compact double locus, and the exact frontier preimage. The marks and their
homotopies survive. It does not yet prove injectivity of the projection.

The resulting finite double graph consists of circle components and
proper arc components. Minimize the number $b$ of components meeting the
source rims among all such essential marked ordinary annuli. If $b>0$,
choose a component meeting the outer rim, reversing the annular depth
coordinate if necessary. If it also meets the inner rim, use a spanning
two-strip tube; otherwise use the nonspanning chain of strip exteriors.
Resolve its two crossing sheets and select the annular successor retaining
an essential rim. In the spanning case there are two candidates. Write
$A,B$ for the two retained rim paths and $E_0,E_1$ for the joining
paths across the old tube end. The old complete rim represents
$A(E_0BE_1^{-1})^{-1}$. The resolving paths in that same end disk
give the two new rim loops. If both contracted, their disk homotopies
would contract this old word. At least one candidate is therefore
essential; its other rim is freely homotopic to it along the annulus.
The nonspanning construction retains the essential rim through its
chain of exterior pieces.
The full tube preimage consists of the two selected strips; this equation
rules out extra intersections with the retained part. The surviving double
relation is the old relation on the retained source copies. Ordinary
crossing charts transport through these copies, while the selected
component disappears. Hence the successor has strictly smaller $b$, a
contradiction.

At $b=0$ every double component is an interior circle. The circle surgery
of Lemma~\ref{lem:dehn-annulus-decrease} treats self-paired, contractible,
and essential paired circles. Its retained-source relation gives a
strict decrease of the finite circle count, with no new boundary double
component. Induction ends at an embedding. Each whole rim is retained
up to its own circle homeomorphism, so its mark and essentiality persist.
Exact pointwise rim parametrizations are restored by the collar
reparametrization when the later marked-cut construction requires them.
Thus three different finite arguments are used: exception repair,
boundary-component minimization, and interior-circle reduction.

For the separate lower-handle application, Theorems
\ref{thm:dehn-protected-disks}, \ref{thm:dehn-protected-annulus}, and
\ref{thm:dehn-enclosing-region} supply the radius-$3/2$ marked disks or
annulus and their enclosing region. The same carrier contains the unit
core, meets the old frontier in the specified patch, and retains a strict
outer transverse margin. Those statements concern protected handle
placement; the rigidity proof below also needs the ordinary tower descent
just established.

\section{Index one}
\label{sec:relative-rigidity-index-one}
\label{sec:index-one}

\subsection{Constructing the relative homeomorphism}

Normalize the lattice by a real linear equivalence taking an integer
basis to the fixed period basis. Pull back both chart families and the
given homotopy through this same equivalence. Thus the old boundary
points, not merely their homotopy classes, remain marked.

In the normalized interval--torus choose two regular values of one
angular coordinate of the comparison map. Relative compression removes
kernels in the phase surfaces while preserving their complete boundary
rims; irreducibility removes the closed spherical components. The
finite component minimum leaves two whole source annuli. Their marked
frontier maps identify the two complementary slabs with the boundaries
of standard solid tori. Use the same annulus maps in both slabs.

Pull a standard meridian band back through each frontier map. Its
meridian is nontrivial in the slab boundary but contracts in the slab:
the ambient fundamental-group injection reflects its standard filling
into this actual slab. The marked proper-disk theorem
of Section~\ref{sec:dehn-proper-disk} supplies an embedded disk.
The collar correction changes its essential band rim to the exact
middle meridian and extends the prescribed parametrization over the
disk. Thicken the disk, retaining both full cap disks and the lateral
annulus. After cutting, the frontier is a sphere. Push this sphere
slightly inward, fill it by source irreducibility, and reattach the
collar to obtain the cut ball with its prescribed complete frontier.
Extend the frontier PL homeomorphism over the ball by coning in finite
ball coordinates. Reglue the two disk faces using their identical
parameters. The two resulting slab maps agree on both entire phase
annuli and fix both old boundary tori. Finite closed subdivisions on
the slabs prove PL regularity across the cuts.

This gives a boundary-fixed PL homeomorphism $E$ from the standard to
the source structure. It may have two angular windings. In coordinates
$[0,1]\times(\mathbb R/p\mathbb Z)^2$, lift the two angular errors from
the lower boundary. On the upper boundary each error is a constant
integer multiple of $p$: it is continuous, lattice-valued, and the
boundary torus is connected. Choose integers $n_1,n_2$ and precompose
with the PL shear
\[
 (s,[y_1],[y_2])\longmapsto
 (s,[y_1+p n_1s],[y_2+p n_2s])
\]
which cancels these errors. It fixes both boundary tori pointwise.
The remaining lifted errors interpolate to zero relative to both ends.
Writing $K$ for the corrected homeomorphism, this gives
$K\simeq\operatorname{id}$ relative to the boundary.
Postcompose that homotopy by $K^{-1}$ to obtain
$\operatorname{id}\simeq K^{-1}$. Set $g=K^{-1}$; finite local inversion
gives both PL directions. Concatenate with the reverse of the original
identity-to-$\phi$ homotopy to obtain $\phi\simeq g$.

\subsection{The supported handle application}

For $H=[-1,1]\times(\mathbb R^2/L)$, the protected annulus and its meridian
cuts give a compact marked ball in the same carrier.  A protected prime
replacement supplies an atlas $e'$ on the same ambient space and an open $N$
containing the
protected ball and the whole frontier, with both identity maps $e\to e'$ and
$e'\to e$ PL on $N$, as in
Theorem~\ref{thm:protected-prime-replacement}. Boundary-proper approximation
supplies the comparison map and its identity homotopy. The rigidity theorem
just proved produces $g$ with both relative homotopies, and
the bounded-lift construction gives a lift $G$ fixing the entire lifted
frontier.

Write $D=\overline B^\iota_\infty\times\mathbb R^\kappa$ for the
lifted cylinder and $C=\overline B^3_\infty(0,1)$ for its compact unit core.
The compactification $p:\mathbb R^3\to B_\infty(0,2)$ is a partial
homeomorphism with full source, PL locally in both directions. It is
the identity on $C$ and on $D\cap\overline B_\infty(0,r)$ for a chosen
$1<r<2$. The protected compact region $P$ lies in this retained set.
The ambient
correction $A$ satisfies
\[
 A(px)=p(Gx)\quad(x\in D),
\]
and is joined to the identity by an isotopy fixed on
$\{\|x\|_\infty\geq2\}\cup D^c\cup\operatorname{fr}D$.
The protected image is literally $A(P)$. Its complementary region in
the fixed box is handled by one proper meridian, a common-width disk
product and a cut ball. Extend the marked frontier comparison over
that ball and reglue; agreement on the full meridian disks, as well as
their rims, is needed. Paste this complementary map to the middle-ball
map along the complete protected annulus and extend by the identity
outside the box. This constructs a standard finite PL placement $Q$
with
\[
 Q(C)\subseteq A(P),\qquad
 B=A^{-1}\circ Q,\qquad B(C)\subseteq P.
\]
To check the inverse domain, let $\pi:D\to H$ be the quotient and set
$q(x)=\pi(p^{-1}Qx)$ for $x\in C$. Placement ensures this formula is
defined. For $y=A^{-1}Qx\in P$, the retained equation $py=y$ gives
$pGy=Ay=Qx$, hence
\[
 g^{-1}q(x)=\pi(A^{-1}Qx)\in\pi(P)\subset N.
\]
This is exactly where the PL identity from $e'$ back to $e$ applies.
The PL inverse of $g$, the standard-chart formula for $q$, and that
identity show that $h\circ B$ is finite PL on $C$, where $h$ is the
original chart to be corrected. No global PL assertion about $B$
in the original Euclidean structure is needed.

Both factors fix the prescribed support, so
\[
 Bx=x\quad\text{if }\|x\|_\infty\geq2
 \text{ or }x\in D^c\cup\operatorname{fr}D.
\]
The supported Alexander construction produces an isotopy $B_t$ from
the identity to $B$ satisfying these same equations at every time.
Concretely for $t>0$ scale the correction as $tB(x/t)$ on its shrinking
support and use the identity outside, with $B_0=\operatorname{id}$.
Bounded support proves continuity at $t=0$; the star-shaped cylinder
and its fixed complement preserve the cylinder-frontier condition.

\section{Index two}
\label{sec:relative-rigidity-index-two}
\label{sec:index-two}

For the solid-torus rigidity theorem, the boundary-fixed comparison
identifies the original meridian band with the standard band. Its middle
circle contracts in the solid torus and is essential in the band.
The marked proper-disk construction gives an embedded source disk
whose rim is essential inside that band. In a finite collar model, the
essential polygon and the required middle circle bound annular regions.
Attach them at distinct collar heights to move the rim to the exact
middle circle, preserving properness and embeddedness; extend its given
boundary parametrization over the disk. Cut along its disk product.
Source irreducibility fills the inward-pushed cut sphere, and collar
reattachment identifies the whole complement with a prism. A PL
extension of its marked boundary map agrees on both copies of the disk,
so regluing the original period yields the relative homeomorphism.
Transport back through lattice normalization gives both PL directions
and both boundary-relative homotopies, as in index one.

The following is the separate placement step in the lower-handle
application. It takes place after the protected comparison $A$ has
been constructed.

For $H=\overline B^2\times(\mathbb R/L)$, the marked middle ball has a side
patch and two disjoint disk faces, each meeting that patch in its prescribed
rim. Rim-fixed PL maps on the two disks
are first pasted with the identity on the side patch and then pasted along the
second rim; the overlap equations give one homeomorphism of the complete
middle-ball boundary.  Two cap balls are then chosen in one fixed box.  Their
contact equations are
\[
 R_{\mathrm{mid}}\cap C_0=D_0,\qquad
 (R_{\mathrm{mid}}\cup C_0)\cap C_1=D_1,\qquad
 (R_{\mathrm{mid}}\cup C_0)\cup C_1=\text{box}.
\]
Take a common finite subdivision of the three PL extensions. They
agree on the whole common disks, and each maps a disk precisely onto
the corresponding target disk. This last equation prevents a point off
an overlap from entering an overlap and proves injectivity of the pasted
map. Denote the source middle ball by $S$ and the target middle ball by
$T=A(P)$. Closed extension to the whole ambient space gives $Q$ with
\[
 Q(S)=T,\qquad Q=\operatorname{id}\text{ on the complement of the box's
 interior}.
\]
The source middle ball $S$ contains $C$. Thus $Q(C)\subset A(P)$.
Set $B=A^{-1}Q$ and apply exactly the inverse-domain calculation above.
It proves finite PL regularity of $hB$ on $C$ and an isotopy fixed on
$\{\|x\|_\infty\geq2\}\cup D^c\cup\operatorname{fr}D$.

\section{Index zero and the lower-handle consumer}
\label{sec:relative-rigidity-index-zero}
\label{sec:index-zero-final}

For $(|\iota|,|\kappa|)=(0,3)$, $H$ is the closed three-torus
$(\mathbb R^3/L)$ and its interval-factor boundary is empty.  The puncture
appearing in the Wall construction is an auxiliary ambient puncture; it is
not the boundary of $H$. First we prove rigidity on the closed torus, then
explain that separate marked-side application.

\subsection{The phase hierarchy and its failure alternative}

Normalize the lattice and write the target as three circles of a common
period $p$. A phase is an inverse image of a regular value of a circle
coordinate of the comparison map. Finite PL position gives finite
surfaces and their two-sided product neighborhoods. Sphere removal and
proper-disk compression give incompressible nonspherical first phases:
sphere components are filled using irreducibility; compression replaces
an essential two-sided annular neighborhood by the two disk faces.
After discarding spheres, the sum $\sum\max(2g-1,0)$ over component
genera strictly decreases under essential compression: separating
compression reduces it by one, nonseparating compression by two for
$g\geq2$, and compression of a torus by one. Thus compression terminates.
Supports avoid the other selected phase and preserve the
complete cut frontier. Among all resulting maps homotopic to the original
map choose one with the smallest sum of the two first-phase component
counts. This is a minimum among actual incompressible phase configurations,
not among unmarked surfaces.

Each closed phase component is an orientable surface with nontrivial
fundamental group injecting into that of the torus. Its group is therefore
abelian. The classification of compact orientable surfaces excludes genus
at least two, and nonsphericity excludes genus zero; hence it is a torus.
The induced tangential map is homotopic to the torus covering determined
by its injective integer matrix. Install these coverings on complete
phase collars, keeping the normal coordinate unchanged. Consequently
the actual phase sets, their counts, and both complementary regions are
unchanged. On either whole frontier the two tangential coordinates of
the installed map are now a covering.

Inside the regions between first phases, cut successively at regular
second and third coordinates, retaining all prior frontiers.
The source surfaces are annuli and disks. Their finite PL normal forms
are installed by homotopies fixed on their complete rims. If the annulus
end labels differ and the disks map homeomorphically, the complementary
cut balls extend the boundary maps and reglue to the desired
homeomorphism. In a folded annulus, the equal end labels give a full
proper interval fiber whose image is constant. A disk failure supplies
the corresponding interval fiber, extended along the installed annulus
collars to the old first frontier. Thus the only other outcome is an
embedded proper arc $k:[0,1]\to R$ with distinct endpoints on
$\operatorname{fr}R$ for which $\psi k$ contracts relative to those
endpoints. Here $\psi$ is still homotopic to the original $\phi$.
These are finite fiber and cut constructions; no assumption that the
annuli are unfolded is inserted.

The first-phase minimum supplies the marked boundary data of every
connected component $P$ of this retained failure region $R$. There
are disjoint whole frontier tori $S_0,S_1$, on which the first circle
coordinate of $\psi$ has the distinct values $a,b$ respectively.
Both tori inject on fundamental groups into $P$, and a path between
their basepoints identifies their image subgroups up to commensurability
(their intersection has finite index in each). The whole component
frontier satisfies $\operatorname{fr}P=P\cap\operatorname{fr}R$.
Here are the two reasons these
assertions hold. First, the normal coordinate takes values in a proper
closed circle arc and can be contracted there. The whole map is
injective on fundamental groups of $P$, so its tangential map is also
injective. Each compact boundary covering has finite-index image in
the tangential torus group, hence each boundary inclusion has
finite-index image in $\pi_1(P)$. Transport along any connecting path
preserves this property and gives commensurability.

Second, a component cannot have its entire frontier at just one phase.
Its frontier is nonempty: the component is a proper closed subset of
the connected ambient torus, since the original identity homotopy
forces surjectivity of the circle coordinate. Restrict the installed
inward collar to this whole frontier. If it had only one phase label,
the normal coordinate on the component and this collar could be pushed
past that label, retaining the other phase and the original map off
the collared region. This removes a nonempty family of first-phase
components without creating any, and the retained incompressibility
data remain admissible. It contradicts the minimum. Thus one can choose
whole boundary components of both phase labels, without asserting that
the endpoints of an arbitrary failure arc lie on different components.

\subsection{A whole product excludes the failure arc}

We use the following standard product criterion in precisely this
situation. A connected orientable irreducible PL three-manifold with
two distinct incompressible boundary tori whose groups are commensurable
after basepoint transport is a product between those tori
\cite[Lemma 5.1, pp.~72--73]{Waldhausen1968}. Indeed finite index makes a
nonzero power of every element of the first torus group lie in the
transported second group, which is the hypothesis of that lemma.

The marked implementation constructs this product in the original
atlas. Two generating circles of $S_0$ give spanning annuli to $S_1$
by the proper annulus argument above. Their upper rims are represented
by finite periodic polygonal windows. Relative general position and
circle/arc removal give the paired annuli and their full product
neighborhoods with the prescribed boundary markings. A regular
neighborhood of the two annuli and the punctured boundary tori has a
remaining sphere boundary; irreducibility fills it. Reattaching the
retained neighborhoods yields the whole product, not just a product
neighborhood of one annulus. The finite realization retains affine
projections recovering the original charts, so the product and its
inverse are PL in those charts. Normalize its lower end to the identity.
Thus there is
\[
 J:S_0\times[0,1]\longrightarrow P,\qquad
 J(x,0)=x,\quad J(S_0\times\{1\})=S_1,
\]
a homeomorphism with the exact equation
\[
 J(x,t)\in\operatorname{fr}R
 \quad\Longleftrightarrow\quad t=0\ \hbox{or}\ t=1.
\]

The connected image of the failure arc lies in the component of $k(0)$.
Write $J^{-1}k(t)=(u(t),v(t))$. Since $\psi k$ contracts relative to
its endpoints, $\psi k(0)=\psi k(1)$. The distinct phase labels $a,b$
force $v(0)=v(1)=d$, where $d\in\{0,1\}$. The explicit homotopy
\[
 (s,t)\longmapsto J\bigl(u(t),(1-s)v(t)+sd\bigr)
\]
fixes both endpoints and moves $k$ to an arc $\beta$ in the original
frontier. Let $c:\operatorname{fr}R\to T^2$ be the installed tangential
covering. Composing this homotopy with $\psi$ and its two tangential
coordinates shows that $c\beta$ contracts relative to its endpoints.
Homotopy lifting for a covering then forces $\beta(0)=\beta(1)$:
the lift of the constant endpoint loop starting at $\beta(0)$ is
constant, and lifts of endpoint-fixed homotopic paths have the same
last point. This contradicts the distinct endpoints of $k$.
The failure alternative is impossible.

This proves the zero-index rigidity theorem. The two other cases
proved above and inverse lattice conjugation prove the stated theorem
for all finite coordinate types and full lattices. Concatenation with
the original identity homotopy supplies the second required homotopy.

\subsection{The marked puncture side in the application}

For the initial zero-handle capping, let $z$ be the auxiliary puncture
in the ambient torus and choose a standard quotient chart $c$ containing
$z$. Set $Y=X\setminus\{z\}$ and let $A\subset Y$ be the compact set to
retain, including the small original quotient cube. Before applying
Theorem~\ref{thm:protected-wall-core} in $Y$, enlarge the protected set
to $A\cup(X\setminus c.\operatorname{source})$. This added set is compact
and misses $z$. Wall gives a compact core $K\subset Y$ with spherical
frontier and with that entire protected set in its interior. Regard $K$
as a subset of $X$. Then
\[
 X\setminus\operatorname{int}_X K\subset c.\operatorname{source},
 \qquad z\in X\setminus K,\qquad A\subset\operatorname{int}_X K.
\]
These equations select the missing side, separate it from the protected
exterior, and put its entire closure in one chart. In chart coordinates
this compact complementary region has exactly the sphere as frontier.
Its membership is constant on each component off that sphere; compactness
excludes the unbounded component, and its nonempty interior selects the
bounded one. Brown applied there
identifies this literal complementary side as a ball with boundary
$\operatorname{fr}_X K$. The retained charts and their old values on
$A$ are kept when the ball is capped.

This uses exactly the compact-core theorem's quantifier over every compact
protected set and its exact frontier equality; it does not require a
stronger Wall theorem. The later zero-index protected ball can be
chosen as a small affine cube inside one retained chart.
Theorem~\ref{thm:protected-prime-replacement} then supplies an irreducible
atlas on the same torus and both PL identity directions on an open
neighborhood of that ball. Its construction is separate from the
closed-torus rigidity proof.

The final lower-handle theorem has three independently quantified rigidity
families, for $(0,3)$, $(1,2)$, and $(2,1)$.  The current consumer first
substitutes these three constructed families into the four-input lower-case
assembly.  That assembly then also substitutes the constructed Brown and
boundary-proper approximation producers.  Its remaining inputs are precisely
the Wall compact cores, the generalized protected-Dehn record, and the three
same-carrier protected prime replacements.  Thus the rigidity families are
actual producers in the current source, including index zero; they are not
premises left unresolved by this chapter.

\section{Sources and cross-chapter contracts}
\label{sec:relative-rigidity-sources}
\label{sec:relative-contracts-limitations}

Hamilton, \emph{On the triangulation of 3-manifolds}, Lemmas 2--3 and the lower
handle application, pp.~64--68, is the primary geometric source
\cite[Lemmas 2--3 and application, pp.~64--68]{Hamilton1976}.
Brown's \emph{A proof of the generalized Schoenflies theorem},
Bull.\ Amer.\ Math.\ Soc.\ 66 (1960), pp.~74--76, supplies Theorems
1--4 on pp.~74--76 and Theorem 5 on p.~76.
His \emph{Locally flat imbeddings of topological manifolds},
Ann.\ of Math.\ (2) 75 (1962), pp.~331--341, supplies the collar
theorem (Theorem 1, p.~337) and bicollar theorem (Theorem 3, p.~339).
Hudson,
\emph{Piecewise Linear Topology}, Lemma 4.2, pp.~92--94, is the standard
relative approximation input. Waldhausen's \emph{On irreducible
3-manifolds which are sufficiently large}, Ann.\ of Math.\ (2) 87
(1968), pp.~56--88, supplies the product criterion, Lemma 5.1,
pp.~72--73. His Theorem 6.1, pp.~77--79, requires boundary
incompressibility \cite[Theorem 6.1, pp.~77--79]{Waldhausen1968}
and is not applied to the solid torus: the
proper-meridian argument above supplies that case directly.

The predecessor inputs are Theorem~\ref{thm:protected-wall-core} on the
punctured or open ambient subspace, Theorem~\ref{thm:protected-prime-replacement}
on the same marked handle, and the proper-disk and protected-surface
theorems cited above. The output is finite PL regularity of the corrected
original chart on its unit core, together with an isotopy fixing both the
entire prescribed cylinder frontier and the radius-two exterior. These
are the data used in the subsequent compatible-smoothing construction.
